\section{Reproducing the direct constructions and the field counterexamples}
\label{sec:p3-reproduction}

All files in this section are in the new
\nolinkurl{research/phase3/} directory. The previous complete paper,
the preceding source files, and the released quartic remain unchanged.
The standalone paper is
\nolinkurl{diophantine_optimization_20261106_v3.tex}.
It contains its own bibliography and does not require external
\LaTeX{} inputs. Static checks verify balanced environments and braces,
unique labels and bibliography keys, and defined references and
citations. No rendered-PDF verification is claimed.

\subsection{Indexed formulas and quartic tables}

The four complete coefficient sources are
\nolinkurl{compression/source-41-direct.json.gz},
\nolinkurl{compression/source-43-direct-semilocal.json.gz},
\nolinkurl{compression/source-41-shift2402.json.gz}, and
\nolinkurl{compression/source-43-semilocal-shift2402.json.gz}.
They are produced by \nolinkurl{compression/build_direct_source.py}.
The options \texttt{--anchor finite} and
\texttt{--anchor semilocal} select the size condition;
\texttt{--beta-shift 2402} selects the translated input.

The generic compiler and checker from the preceding construction
are reused unchanged. From the repository root, one example is
\begin{verbatim}
python3 -E -B research/phase2/finite_fibre/compile_shared_source.py \
  research/phase3/compression/source-41-direct.json.gz \
  --output research/phase3/finite_fibre/polynomial-41
python3 -E -B research/phase2/finite_fibre/verify_shared_specification.py \
  research/phase3/compression/source-41-direct.json.gz \
  research/phase3/finite_fibre/polynomial-41/polynomial-specification.json
\end{verbatim}
The other three sources have separate matching specification packages
and passing receipts. The verification concerns integer coefficient
conversion, triangular monicity, guards, all exceptional quadratic
cases, the finite index sets, and degree and coefficient majorants.
It is not a verification of the external arithmetic hypotheses.

The direct quartic producer is
\nolinkurl{quartic/build_direct_quartic.py}. Reproduction uses
\begin{verbatim}
python3 -E -B research/phase3/quartic/build_direct_quartic.py \
  --share-height --precision 52 --scalar-order 32 \
  --output research/phase3/quartic/quartic-direct
python3 -E -B research/phase2/quartic/verify_shared_quartic.py \
  research/phase3/quartic/quartic-direct/polynomial.json.gz
\end{verbatim}
Add \texttt{--small-coefficients} and use the corresponding
small-coefficient output directory for the 1339-variable polynomial.
The checker imports no producer. It checks every gate, the reversible
row operations and eliminations, and the complete polynomial obtained
by independent exact squaring and convolution of the final residuals.
No floating-point or randomized comparison is used.

The preferred tables additionally use the certified transformations in
\nolinkurl{quartic/transform_direct_quartics.py}. The corresponding
independent checker is
\nolinkurl{quartic/verify_final_transforms_independent.py}.
It checks every affine and constant substitution, the inverse coordinate
maps or unique lifts, and both the complete substituted polynomial and
the independently squared residual list. It imports no producer. The
fresh transformed packages preserve their frozen parents for comparison.
The distinction between encoded coefficient entries and coordinate
monomials after specializing a fixed integer is checked separately by
\nolinkurl{compression/check_fixed_coefficient_support.py}.

\subsection{Independent checks for the field assertions}

The following scripts use exact arithmetic and have separate receipts:
\begin{description}
\item[Quartic base change.]
\nolinkurl{number_fields/check_quartic_base_change_20261106_v1.py}
checks the generic $\mathbb Q(i)$ section against both the complete
tables and the independently squared restricted residuals, for all
four quartics.
\item[Rejected contacts.]
\nolinkurl{number_fields/check_contact_exclusions_20261106_v1.py}
checks the three exact input relations in the four unshifted indexed
sources and four quadratic systems. The shifted sources are expressly
outside this counterexample's scope.
\item[The totally real extension.]
\nolinkurl{number_fields/check_real_anchor_20261106_v1.py}
checks the field, point, and doubling calculations in the displayed
degree-eight compositum.
\item[Division norms.]
\nolinkurl{number_fields/build_division_norm_specs_20261106_v1.py}
produces all four determinant specifications. The independently
written
\nolinkurl{number_fields/verify_division_norm_specs_independent_20261106_v1.py}
imports neither this producer nor an earlier checker. It compares
every embedded residual with pinned source bytes, reconstructs each
matrix-entry index, verifies formal Galois conjugacy, checks the
explicit inert primes, and proves the stated exact degrees from the
unit coefficient of the highest pure input power. It also recomputes
the coefficient bounds and the unrounded finite-algebra degree
majorants before their displayed powers-of-two rounding.
\end{description}

The accompanying independent proof review explains the local
division arguments, including the degree-36 norm over
$\mathbb Q(i)$, and the coprime-extension module argument. The norm
determinants are exact finite specifications, not expanded tables.
The counterexamples and norm equivalences are unconditional algebraic
results. They do not certify that the changed source defines any
extension field's ring of integers.

The final affine and constant-chain transformations have separate
number-field certificates and degree-72 specifications with dated
suffix \nolinkurl{20261106_v2}. These are checked independently in
addition to the preceding specifications. In particular, elimination of
a defining residual preserves common zero sets over every field, but
does not by itself preserve every cancellation among nonzero squares
over a nonreal field. The explicit Gaussian sections are therefore
verified again for each new literal quartic.

\subsection{Rebuild and integrity records}

\nolinkurl{assemble_phase3.py} builds this paper from the preceding
complete report and the new mathematical supplements, and checks that
the preceding report's bytes are unchanged. The verification index
records the passing receipts and hashes of the checked artifacts.
\nolinkurl{make_checkpoint.py} packages the complete research tree,
the pinned source manuscript, and the unchanged released baseline.
It hashes the same bytes it writes and verifies the resulting ZIP
archive. Its manifest makes the checkpoint independently inspectable.
No commit, push, or pull-request submission is part of these scripts.
